% from ega3-1-fr.tex
\oldpage[III]{110}
on a
\[
\chi_A(\mathcal F(n))=\operatorname{long}H^0(X,\mathcal F(n))
=\operatorname{long}\Gamma(X,\mathcal F(n))
\]
pour $n$ assez grand, d'où (ii)~; cela entraîne
$\chi_A(\mathcal F(n))\geq0$ pour $n$ assez grand, et (iii) résulte donc
de (i)~; comme d'ailleurs l'assertion d'unicité de (i) est immédiate, il
reste à prouver l'existence du polynôme $P$.

Montrons d'abord qu'on peut supposer que $\mathfrak m\mathcal F=0$, où
$\mathfrak m$ est l'idéal maximal de $A$. En effet, il existe un entier
$s>0$ tel que $\mathfrak m^s=0$, et $\mathcal F(n)$ admet donc une
filtration finie
\[
\mathcal F(n)\supset\mathfrak m\mathcal F(n)\supset\cdots
\supset\mathfrak m^{s-1}\mathcal F(n)\supset0.
\]
Par récurrence, on déduit de
(\hyperref[III.2.5.2.1-fr]{2.5.2.1}) que
\[
\chi_A(\mathcal F(n))=\sum_{k=1}^{s}
\chi_A(\mathfrak m^{k-1}\mathcal F(n)/\mathfrak m^k\mathcal F(n))~;
\]
comme
$\mathfrak m^{k-1}\mathcal F(n)/\mathfrak m^k\mathcal F(n)
=\mathcal F'_k(n)$, où
$\mathcal F'_k=\mathfrak m^{k-1}\mathcal F/\mathfrak m^k\mathcal F$,
cela prouve notre assertion.

Supposons donc $\mathfrak m\mathcal F=0$~; si $X'$ est le sous-schéma
fermé de $X$, image réciproque par le morphisme structural
$X\to\operatorname{Spec}(A)$ de l'unique point fermé de
$\operatorname{Spec}(A)$, et $j:X'\to X$ l'injection canonique, on a
$\mathcal F=j_*(\mathcal F')$, où $\mathcal F'$ est un
$\mathcal O_{X'}$-Module cohérent~; $X'$ est un schéma projectif sur
$\operatorname{Spec}(K)$, où $K=A/\mathfrak m$. Si
$\mathcal L'=j^*(\mathcal L)$, $\mathcal L'$ est très ample relativement
à $\operatorname{Spec}(K)$ (II, 4.4.10), et on a
$\mathcal F(n)=j_*(\mathcal F'(n))$, où
$\mathcal F'(n)=\mathcal F'\otimes_{\mathcal O_{X'}}
\mathcal L'^{\otimes n}$ (0\textsubscript{I}, 5.4.10). On en conclut que
$\chi_A(\mathcal F(n))=\chi_K(\mathcal F'(n))$ (G, II, 4.9.1), et on est
donc ramené au cas où $A$ est un corps.

Notons maintenant que $X$ peut être considéré comme un sous-schéma
fermé de $P=\mathbf P_A^r$ pour un $r$ convenable (II, 5.5.4, (ii))~;
si $i:X\to P$ est l'injection canonique, on voit comme ci-dessus que l'on
a $\chi_A(\mathcal F(n))=\chi_A(i_*(\mathcal F)(n))$, de sorte que l'on
peut se borner au cas où $X=\mathbf P_A^r=\operatorname{Proj}(S)$ avec
$S=A[T_0,\ldots,T_r]$, $A$ étant un corps.

Cela étant, on a $\mathcal F=\widetilde M$, où $M$ est un $S$-module
gradué de type fini (II, 2.7.8)~; il existe par suite une résolution finie
de $M$ par des $S$-modules gradués libres de type fini
\[
0\to L_q\to L_{q-1}\to\cdots\to L_1\to M\to0
\]
en vertu du th. des syzygies de Hilbert (M, VIII, 6.5)~; comme
$M\mapsto\widetilde M$ est un foncteur exact en $M$ (II, 2.5.4), on a
aussi une suite exacte
\[
0\to\widetilde L_q\to\widetilde L_{q-1}\to\cdots
\to\widetilde L_1\to\widetilde M\to0
\]
et par suite, pour tout $n\in\mathbb Z$, la suite
\[
0\to\widetilde L_q(n)\to\widetilde L_{q-1}(n)\to\cdots
\to\widetilde L_1(n)\to\widetilde M(n)\to0
\]
est exacte~; appliquant par récurrence sur $q$ la prop. (2.5.1), il vient
\[
\chi_A(\widetilde M(n))=\sum_{j=1}^{q}(-1)^{j+1}
\chi_A(\widetilde L_j(n))
\]
et pour prouver (i), on est donc ramené au cas où $M$ est libre et gradué
de type fini, donc au cas où $M=S(h)$ pour un $h\in\mathbb Z$. Comme
alors $\widetilde M(n)=(M(n))^\sim=(S(n+h))^\sim$ (II, 2.5.15), on voit
finalement que le théorème résultera du
\end{proof}

